% Pursuing Stacks (cotes 134-2 to 134-8): archive cote 134-4, site page 89 (PDF page 90), typescript page 176.
% typescript. Draft compared: ps3.tex:4558-4637.
% AI-transcribed by Claude Opus 5.5 (claude-opus-5-5, adaptive thinking, effort high), Message Batches API, batch ps2 (msgbatch_01NVwr1Yuo8qJf5rykeGBXPK), 2026-11-12; not yet checked by eye.
\scanpage{134-4/89}
\mspage{176}

a$_2$) The functor $\overline{i^*}$ induced by $i^*$ on the localizations (which exist\unsure{s}
because of a$_1$))
\[ W^{-1}M \longrightarrow (\mathrm{H}\over{e}{o}\mathrm{t}_A) \overset{\hand{def}}{=} \underline{W}_A{}^{-1}A^{\wedge} \]
is an equivalence.

\uline{Definition} 6. The functor i is called a \inserted{W-}\uline{test} \uline{functor} if this functor
\uline{and} the induced functors $i_{/a} : A_{/a} \longrightarrow M$ (for a in A) are weak $\underline{W}$-test
functors.

In view of the definition 3, this amounts to the two conditions:
a) A is a $\underline{W}$-test category , \hand{i.e. the functors $i^*_A$ and $i^*_{A_{/a}}$ are model preserving,} \inserted{(and A is $\underline{W}$-aspheric,)}
b) the functors $i^*$ and$(i_{/a})^*$ from M into the categories $A^{\wedge}$ and
$(A_{/a})^{\wedge}$\hand{$\xrightarrow{\approx}$}$A^{\wedge}{}_{/a}$ are model preserving (for W and $\underline{W}_A$ resp. $\underline{W}_{A_{/a}}$ ) .

\uline{Example} Consider the canonical functor induced by $i_A$ \hand{$: A^{\wedge} \to (\mathrm{Cat})$,}
\[ i^{\mathrm{o}}_A : A \longrightarrow (\mathrm{Cat}) \quad , \quad (\mathrm{Cat}) \text{ endowed with } \underline{W} \ , \]
this functor is a weak $\underline{W}$-test functor (resp. a $\underline{W}$-test functor) iff A is
a weak $\underline{W}$-test category (resp. a $\underline{W}$-test category).

These two definitions are pretty formal indeed. Their justificati-%
on is mainly in the two theorems below.

\xout{Theorem 1}\xout{xxx} \uline{We assume} \xout{that W} \uline{from now on that the basic localizer} $\underline{W}$
\uline{satisfies the fibration axiom} L 5) of page 164. Also, we recall
that an\xout{\illegible{1}} object X in (Cat) is \uline{contractible} (for the canonical
homotopy structure of (Cat)) iff X is non-empty and the category
\uline{Hom}(X,X) is connected - indeed it is enough even that id$_X$ belong to the
same connected component as some \uline{constant} map from X into itself. This
\margin{\unsure{Sa généralisation p. 252}}condition is satisfied for instance if X has a\xout{n} final or an initial
object.

\uline{Theorem} \hand{1}. We assume that M = (Cat), W = $\underline{W}$ i.e. we've got a functor
\[ i: A \longrightarrow (\mathrm{Cat}) \quad , \quad (\mathrm{Cat}) \text{ endowed with } \underline{W} \ , \]
and we assume that for any a in A, i(a) is contractible (cf. above), i.e.
that i factors through the full subcategory $(\mathrm{Cat})_{\mathrm{cont}}$ of contractible
objects of (Cat). The following conditions are equivalent:
(i) i is a \inserted{$\underline{W}$-}test functor (def. 6).

(i') For any a in A, the induced functor $i_{/a} : A_{/a} \longrightarrow (\mathrm{Cat})$ is a
weak \inserted{$\underline{W}$-}test functor, and moreover A is $\underline{W}$-aspheric .

(ii) $i^*(\hand{$\Delta$}_1)$ is $\underline{W}_A$-aspheric over $e_{A^{\wedge}}$ i.e. the products $axi^*(\hand{$\Delta$}_1)$
in $A^{\wedge}$ are $\underline{W}$-aspheric, for any a in A\hand{, and A is $\underline{W}$-aspheric.}

\margin{89}
